\subsection{Link-quality-aware routing}
FANETs are characterised by high mobility and rapidly changing
topology~\cite{bekmezci}, and both the general UAV network
surveys~\cite{gupta,hayat} and the routing-specific
surveys~\cite{arafat,lakew} treat link durability as a decisive factor.
OLSR~\cite{rfc3626} selects routes by hop count: a link counts as usable as
soon as HELLO messages traverse it in both directions, and its packet loss does
not enter the path cost. OLSRv2~\cite{rfc7181} adds a link-metric framework, but
the OLSR implementation in ns-3 follows RFC 3626. ETX~\cite{decouto} and
ETT~\cite{draves} replace hop count with a cost reflecting loss and bandwidth;
olsrd~\cite{olsrd} brings ETX into OLSR and is the natural distributed baseline
for any link-quality-aware routing method.

\subsection{Link-quality forecasting and GNNs}
Link-quality estimation is a research direction in its own right, ranging from
packet-statistics estimators~\cite{baccour} to machine-learning
models~\cite{cerar}. GNNs learn node representations through iterated message
passing: GCN~\cite{kipf}, GraphSAGE~\cite{graphsage} with neighbour aggregation
that generalises to graphs unseen during training, and GAT~\cite{gat} with
attention weights per neighbour. A UAV network is naturally a time-varying
graph, which makes GNNs a reasonable candidate. In wired networks,
RouteNet~\cite{routenet} uses a GNN to predict delay and loss from topology and
routing. In UAV networks, Bai et al.~\cite{bai2025} predict links with a dynamic
GNN and take GraphSAGE and node2vec themselves as baselines;
GNNPPOR~\cite{gnnppor} couples a GNN with PPO for path selection in FANETs; and
a recent review~\cite{drones2026} surveys GNN applications across UAV systems.
Within FANETs, QMR~\cite{qmr} uses Q-learning for multi-objective next-hop
selection, and XGBoost~\cite{xgboost} is the strong tabular baseline that any
GNN evaluation needs. Predictive-OLSR~\cite{rosati2013} uses GPS and relative
velocity to adjust ETX weights, while OLSR-LCN~\cite{liu2025} combines link
lifetime, interference and node load. In other words, \emph{using forecasting,
position information or GNNs for UAV routing is no longer a new idea}; the open
question is what such a model is actually worth when measured by service.

\subsection{The gap between accuracy and service benefit}
Cerar et al.~\cite{cerar} observe in their survey that most link-quality
estimation models are evaluated offline by accuracy and are rarely placed
inside a routing loop to measure the service benefit. Shi et al.~\cite{shi2024}
show that several families of machine-learning models converge on the same
accuracy ceiling imposed by channel stochasticity. In other fields this gap is
well documented: in autonomous driving, the ranking of motion predictors shifts
substantially when evaluation moves from open loop to closed
loop~\cite{closedloop2025}; in information retrieval, much of the reported gain
of neural models disappears once baselines are properly tuned~\cite{lin2018};
and in machine learning for partial differential equations, 79\% of surveyed
articles compare against a weak baseline~\cite{mcgreivy2024}.
Pre-registration~\cite{nosek} is a direct countermeasure against reselecting a
configuration after seeing the results.

\subsection{Relay placement}
Work on relay trajectory optimisation~\cite{zeng2016} controls UAV positions to
maximise throughput, usually with an analytic channel model and the assumption
that traffic follows the relay chain the controller selects. Software-defined
UAV networks~\cite{abir} employ a controller with a global view; the drawbacks
usually cited are signalling overhead, control delay and scalability. This work
does not solve the optimal trajectory problem; it keeps a fixed equispacing
placement rule in order to isolate the \emph{trigger timing}, and additionally
measures a geometric bound so that the distance between that rule and the best
achievable outcome is known.

\subsection{The targeted gap}
To the best of our knowledge, no prior work in UAV networking (i) drives two
different actions from one weight-locked forecaster inside the same simulator,
(ii) pre-registers the endpoints and the test order and then reports
unfavourable results as well, and (iii) contrasts the ranking by forecast error
with the ranking by service on the same set of traces. Those are the three
points this paper targets.
